% --- Back cover ---
\clearpage
\thispagestyle{empty}
\begin{tikzpicture}[remember picture, overlay]
  % --- Wordmark ---
  \node[anchor=north, font={\sffamily\bfseries\fontsize{48pt}{52pt}\selectfont}]
    at ([yshift=-26mm]current page.north) {PERSONIX};
  \node[anchor=north, font=\rmfamily\itshape\large]
    at ([yshift=-44mm]current page.north) {Nekompromisní změna};
  \node[anchor=north, font=\rmfamily\small,
        text width=\paperwidth-50mm, align=center]
    at ([yshift=-54mm]current page.north)
    {Necenzurovatelná a nezkorumpovatelná decentralizovaná reputační síť};
  % --- Body ---
  \node[anchor=north, text width=\paperwidth-50mm, align=justify, font=\small]
    at ([yshift=-72mm]current page.north)
    {%
      Stát fungoval, když komunikace byla pomalá, rozhodnutí lokální
      a autorita žila ve stejném městě, kterému vládla. Dnešní sítě obracejí
      všechny tři předpoklady --- a instituce na nich postavené už nesedí
      do světa, který spravují. Tato kniha popisuje nástroj, který sedí:
      decentralizovanou reputační síť, ve které je identita vaše, pravda
      veřejně auditovatelná a podplacení znamená reputační sebevraždu.

      \smallskip
      Není to manifest. Není to předpověď. Je to funkční plán toho, jak
      lze nástupce státu postavit --- dobrovolně, jednu deklaraci po druhé.
    };
  % --- Footer ---
  \node[anchor=south, font=\sffamily\footnotesize,
        text width=\paperwidth-50mm, align=center]
    at ([yshift=12mm]current page.south)
    {%
      Pavel Kudrna\quad\textbullet\quad Praha, 2026\par
      \href{https://personix.org}{personix.org}\par
      Licencováno pod Creative Commons Uveďte původ-Zachovejte licenci 4.0
      Mezinárodní (CC BY-SA 4.0).
    };
\end{tikzpicture}%
\mbox{}%
\clearpage
